\documentclass[10pt]{article}
\providecommand{\FIGDIR}{fig}
\newcommand{\DIGESTDATE}{4 November 2026}
\input{preamble}

\begin{document}

% ==================================================================== MASTHEAD
\thispagestyle{empty}
\begin{tikzpicture}[remember picture, overlay]
  \fill[Deep] (current page.north west) rectangle
      ($(current page.north east)+(0,-67mm)$);
  \fill[Amber] ($(current page.north west)+(0,-67mm)$) rectangle
      ($(current page.north east)+(0,-68.4mm)$);
  \foreach \x/\y/\r in {22/12/0.9, 41/26/1.5, 63/9/0.7, 88/31/1.1, 112/17/1.9,
                        138/29/0.8, 159/13/1.3, 178/54/1.0, 196/21/1.6,
                        30/44/0.6, 74/57/1.2, 125/48/0.9, 168/6/0.8, 191/59/0.7,
                        52/18/1.0, 100/60/0.6, 148/20/1.1, 66/49/1.4, 205/45/0.8}
    \fill[Sky, opacity=0.30] ($(current page.north west)+(\x mm,-\y mm)$) circle (\r mm);
\end{tikzpicture}

\vspace*{4mm}
{\sffamily\fontsize{7.6}{9}\selectfont\color{Sky}%
 AUTOMATED DAILY LITERATURE DIGEST\;\textbullet\;PhD RESEARCH WATCH\par}
\vspace{2.2mm}
{\sffamily\bfseries\fontsize{27}{29}\selectfont\color{white}%
 Particulate Matter\par}
{\sffamily\fontsize{27}{29}\selectfont\color{Sky}%
 Monitoring \& Health Effects\par}

\vspace{5.0mm}
{\sffamily\fontsize{8.4}{10}\selectfont\color{white}%
 Issue 2026-11-04\;\;\textcolor{Amber}{\textbar}\;\;Entry window 4 November 2026 (full day)\;\;%
 \textcolor{Amber}{\textbar}\;\;18 records screened in scope\par}
\vspace{18mm}

\noindent
\begin{minipage}[t]{0.190\linewidth}\metric{18}{RECORDS IN SCOPE\\(70 EXCLUDED)}{Deep}\end{minipage}\hfill
\begin{minipage}[t]{0.190\linewidth}\metric{7}{OF 10 SUBTOPIC\\BANDS POPULATED}{Teal}\end{minipage}\hfill
\begin{minipage}[t]{0.190\linewidth}\metric{15}{EFFECT ESTIMATES\\WITH A CI}{Sky}\end{minipage}\hfill
\begin{minipage}[t]{0.190\linewidth}\metric{3}{TIER-A\\RECORDS}{Amber}\end{minipage}\hfill
\begin{minipage}[t]{0.190\linewidth}\metric{3}{RECORDS CARRIED ON\\METADATA ALONE}{Clay}\end{minipage}

\vspace{6mm}

\section*{\sffamily\bfseries\color{Deep}Signal of the day}
\vspace{-1.5mm}{\color{Amber}\rule{\linewidth}{1.1pt}}\vspace{2.5mm}

\begin{keybox}
{\sffamily\bfseries\small\color{Deep}A metrology-chain network calibration worked for temperature, humidity and pressure in Helsinki and Antwerp and failed first on PM: feasible only on restricted sensor subsets, because PM signals are localized and weakly correlated.}\\[1.4mm]
\footnotesize
VTT-led; Laplace-domain pairwise transfer functions, a separated spatial term and an overdetermined solve against one reference, with Monte Carlo uncertainty. Helsinki linked to SMEAR III; Antwerp from SensEURCity.
\textbf{Where it is weak:} the abstract gives no PM R$^2$, RMSE or uncertainty, so the size of the failure cannot be graded, and two cities do not bound the geometries where it holds.
\textbf{Why it matters here:} node-to-node calibration transfer assumes neighbouring sensors see related signals. The more local the source field, the more value a dense network has and the less that assumption holds; the property that justifies deployment is the one that defeats reference-free transfer. Co-location does not go away.\par
\end{keybox}

\vspace{3mm}

\section*{\sffamily\bfseries\color{Deep}What changed today}
\vspace{-1.5mm}{\color{Amber}\rule{\linewidth}{1.1pt}}\vspace{2.5mm}

\footnotesize
\begin{itemize}

\item \textbf{Life-course composition, on a 1930--2010 CTM.} UK Biobank, 47,737 with birth-to-baseline residences, 13--14 years of follow-up. Per IQR of PM$_{2.5}$: COPD HR 1.138 (1.069--1.211), asthma 1.171 (1.105--1.241), lower respiratory infection 1.046 (1.002--1.091), lung cancer 1.193 (1.064--1.337); SO$_4$ and EC positive for all four. IQR not in the abstract; most of the period predates any validating monitor.

\item \textbf{A joint HR that is not an interaction.} CFPS China, 9,265 adults: PM$_{2.5}$ HR 1.54 (1.42--1.67) per unstated IQR for incident CVD; high PM$_{2.5}$ with compound dry-hot extremes 1.92 (1.67--2.21), called synergistic with no interaction scale given. An IQR estimate that large sits well above the cohort literature; the abstract names neither exposure source nor CVD ascertainment.

\item \textbf{Tenfold regional PMF in a regulated industry.} NIOSH surveillance, 30,319 underground coal miners, 1991--2024: Central Appalachia PMF PR 10 (5.7--17.5) and CWP PR 3.9 (3.4--4.5), tenure-adjusted; continuous miner operators PMF PR 3.0 (2.3--4.0). Job and region proxy for dust; nothing was measured.

\item \textbf{Sample count as the exposure instrument.} Colorado State: five people ran 83 AirPen personal monitors across one shift at a plant of about 100 workers; 67 valid PM and noise samples, 22 formaldehyde. Intervals grew several-fold at smaller $n$; one shift says nothing about day-to-day variance. \textsc{backfill}, entry 28 May.

\item \textbf{Two sensor evaluations, two standards of evidence.} ARMIE (SPS30, Aarhus): r 0.93--0.98 against a DustTrak but 0.74--0.94 against SMPS+APS, tested leave-one-source-out, the design epidemiology needs. Goyal (two PA-II, one FEM, 150 days): +17.55\,\% raw bias to 0.13\,\% NMB, R$^2$ 0.9010, with no stated train/test split. Both \textsc{backfill}; only the first informs transfer.

\item \textbf{Exposure basis (x1).} Of 18 records, 3 rest on a low-cost sensor or network (network calibration, ARMIE, Goyal) and 1 on a personal monitor (AirPen); none of these is a health study. Health side with a recoverable basis (n=5): 60\,\% measured, 20\,\% modelled, 20\,\% proxy, but that measured share is the AirPen plant survey, six sawmill extracts and a chicken-lung eluate. \emph{No effect estimate below is tied to a measured PM value}: a CTM (Biobank), an unstated surface (CFPS), a pooled mix (rhinitis), job or region (miners).

\end{itemize}

\vspace{2mm}

\begin{keybox}
{\sffamily\bfseries\footnotesize\color{Deep}Window continuity}\\[1.2mm]
{\fontsize{7.2}{8.8}\selectfont
The 3 November issue closed with \texttt{last\_entry\_date} at 3 November. This issue was built on the \textbf{6 November 00:33 UTC (5 November 18:33 CST) run}, a slot that fired 19.5 hours late. Before 22:00 the newest buildable day is yesterday, so \textbf{4 November is covered in full and 5 November waits for the next run after 22:00 CST}. Three \textsc{backfill} records, the cap. No rollup is owed: W45 (1--7 November) falls on Saturday 7 November. \textbf{The OpenAlex connector was not available on this session}, so the provisional 3 November recovery of that leg remains unconfirmed.\par}
\end{keybox}

\vspace{3mm}
\par\ifdim\dimexpr\pagegoal-\pagetotal\relax<18\baselineskip\newpage\fi
\section*{\sffamily\bfseries\color{Deep}Corpus analytics}
\vspace{-1.5mm}{\color{Amber}\rule{\linewidth}{1.1pt}}\vspace{2.5mm}

\noindent\begin{minipage}{\linewidth}\includegraphics[width=\linewidth]{\FIGDIR/f1_subtopics.png}
\figcap{Subtopic assignment of the 18 records across seven of ten bands: sensing and occupational four each, respiratory and exposure assessment three each, mechanistic two, cardiovascular and neuro one each. Reproductive, burden and other clinical are empty today.}\end{minipage}

\vspace{2mm}
\noindent\begin{minipage}{\linewidth}\includegraphics[width=\linewidth]{\FIGDIR/f2_design_endpoint.png}
\figcap{Left --- architecture: measurement campaign six, review or synthesis four, metadata-only three, cohort two, experimental or analytical two, cross-sectional one. Right --- endpoints are singletons apart from the three unrecoverable metadata-only records; many are metrological rather than clinical, which is the day's balance.}\end{minipage}

\vspace{2mm}
\noindent\begin{minipage}{\linewidth}\includegraphics[width=\linewidth]{\FIGDIR/f3_metric_geography.png}
\figcap{Left --- particle metric: unspeciated or emissions 9, PM$_{2.5}$ only 4, PM$_{2.5}$+PM$_{10}$ jointly 4, composition 1. Right --- geography: North America 4, Europe 4 (including the UK), global or multi-region 4 (three reviews, plus one metadata-only record with no location that falls there by default), China 2, and one each Southeast Asia, East Asia ex-China, Middle East and N. Africa (Istanbul) and polar (Ardley Island).}\end{minipage}

\vspace{2mm}
\noindent\begin{minipage}{\linewidth}\includegraphics[width=\linewidth]{\FIGDIR/f6_lifecourse.png}
\figcap{Life-course windows: childhood (rhinitis meta-analysis, trajectory review), adolescence (ABCD), working age (four occupational and indoor settings), and adult to older (UK Biobank, CFPS). No preconception or in-utero record today.}\end{minipage}

\vspace{2mm}
\noindent\begin{minipage}{\linewidth}\includegraphics[width=\linewidth]{\FIGDIR/f4_heatmap.png}
\figcap{Subtopic $\times$ study architecture. Sensing sits in measurement (3) plus one metadata-only; occupational spreads over measurement (2), analytical chemistry and the miner surveillance record. The only cohorts are respiratory and cardiovascular, and neither carries a measured exposure; no measured-exposure cohort is in the matrix today.}\end{minipage}

\vspace{2mm}

\IfFileExists{\FIGDIR/extra_figs.tex}{\input{\FIGDIR/extra_figs.tex}}{}



\noindent\begin{minipage}{\linewidth}\includegraphics[width=\linewidth]{\FIGDIR/f5_forest.png}
\figcap{Ratio scale, 15 estimates from 4 sources: UK Biobank (4), CFPS (3), one meta-analysis (2), NIOSH miners (6). The rhinitis pair carries no increment in its abstract and converts to no concentration contrast. Estimates within a source share a cohort and are not independent. Increments differ (per unstated IQR, categorical contrast, none in the abstract) and PR on prevalence is not HR on incidence, so nothing here is poolable; no I$^2$ is reported. Upper rows are job and region contrasts in respirable coal dust, not ambient PM.}\end{minipage}

\vspace{4mm}

\par\vskip 0pt plus 24\baselineskip\penalty-200\vskip 0pt plus -24\baselineskip
\section*{\sffamily\bfseries\color{Deep}Paper-level digest}
\vspace{-1.5mm}{\color{Amber}\rule{\linewidth}{1.1pt}}\vspace{2.5mm}

{\fontsize{7.4}{9}\selectfont\color{Slate}Ordered by cluster. Each entry gives the
methodological core and the finding that matters, not an abstract paraphrase. Records
marked \textit{metadata only} carry no recoverable abstract and assert no finding.\par}

\band{Neuro / mental health\hfill 1 record}{Deep}

\paperentry{Cardenas-Iniguez ABCD linked external data environment review first decade 2026}
{Environmental influences using linked external data in ABCD: A review of the first decade.}
{Dev Cogn Neurosci}
{Ten years of Linked External Data work in the Adolescent Brain Cognitive Development Study (\textasciitilde{}11,800 US youth), covering neighbourhood conditions, built environment, air pollution, greenspace and state policy. The authors conclude that development reflects cumulative influences and that timing and study design bound what can be concluded, and they call for refined geocoding and exposure methods. No effect sizes in the abstract. Limitation: an internal working-group review of its own dataset, and residential geocoding ignores school and transit time, which matters at this age. Relevance: moderate - the largest US neurodevelopment cohort names exposure resolution as a limiting method.}
{10.1016/j.dcn.2026.101822}

\par\vskip 0pt plus 10\baselineskip\penalty-200\vskip 0pt plus -10\baselineskip
\band{Cardiovascular \& metabolic\hfill 1 record}{Teal}

\paperentry{Shi CFPS compound dry-hot extremes PM$_{2.5}$ joint incident CVD China 2026}
{Joint effects of air pollution and compound dry-hot extremes on cardiovascular disease risk across adults: a cohort study in China.}
{BMC Public Health}
{China Family Panel Studies, 9,265 adults with 69,297 repeated observations, 2010-2018, mixed-effects Cox. Per IQR, PM$_{2.5}$ HR 1.54 (1.42-1.67) for CVD and 1.75 (1.46-2.11) for coronary disease; dry-hot cumulative days HR 1.30 (1.20-1.41). Concurrent high exposure to both gave a joint HR of 1.92 (1.67-2.21), called synergistic, and repeated across PM$_{1}$, PM$_{10}$, NO$_2$ and O$_3$. Limitation: an IQR HR of 1.54 is far above the cohort literature and neither the IQR nor the exposure source is stated; the abstract does not specify how CVD was ascertained; joint exposure and synergy are different quantities, and no interaction scale is given. Relevance: low - no measurement content, but a compound-exposure design that a network could sharpen.}
{10.1186/s12889-026-28927-4}

\par\vskip 0pt plus 10\baselineskip\penalty-200\vskip 0pt plus -10\baselineskip
\band{Respiratory \& allergic\hfill 3 records}{Sky}

\paperentry{Yang UK Biobank life-course PM$_{2.5}$ constituents COPD asthma LRTI lung cancer 2026}
{Life-Course Exposure to PM(2.5) Constituents and Major Respiratory Diseases: Constituent-Specific Associations and Candidate Immune-Metabolic Signatures in a Large Prospective Cohort.}
{Lung}
{47,737 UK Biobank participants with birth-to-baseline residential histories, 13-14 years of follow-up. Per IQR of life-course PM$_{2.5}$: COPD HR 1.138 (1.069-1.211), asthma 1.171 (1.105-1.241), lower respiratory infection 1.046 (1.002-1.091), lung cancer 1.193 (1.064-1.337). SO4 and EC were positive for all four outcomes; mixture associations held for COPD and asthma. The IQR in $\mu$g\,m$^{-3}$ is not given in the abstract. Limitation: decades of exposure rest on one retrospective CTM with no available monitor validation for most of the period, and the 47,737 with histories are a selected fraction of the Biobank. Relevance: moderate - composition, not mass, is the signal, which a mass-calibrated optical sensor cannot resolve.}
{10.1007/s00408-026-00948-1}

\paperentry{Zhu meta-analysis air pollutants allergic rhinitis children PM$_{2.5}$ TRAP mold 2026}
{Association between exposure to gaseous air pollutants and the risk of allergic rhinitis in children and adolescents: A meta-analysis.}
{Int J Pediatr Otorhinolaryngol}
{Meta-analysis of 21 studies (13 cross-sectional, 6 cohort, 2 case-control). Pooled ORs: mold 1.45 (1.16-1.80), traffic-related pollution 1.39 (1.21-1.59), PM$_{2.5}$ 1.23 (1.05-1.43), O$_3$ 1.15 (1.05-1.26), NO$_2$ 1.09 (1.00-1.17); PM$_{10}$, SO$_2$ and CO not significant. Estimates for PM$_{10}$, NO$_2$ and O$_3$ were higher in Chinese studies. Limitation: the abstract gives no exposure increment, so the pooled ratios cannot be compared across pollutants or to any concentration, and 13 of 21 studies are cross-sectional. Mold is not an ambient pollutant, and a 'gaseous' title over a PM$_{2.5}$ estimate shows loose scoping. Relevance: low - a pooled OR with no increment cannot feed a sensor-based exposure-response.}
{10.1016/j.ijporl.2026.113068}

\paperentry{Drummond mini-review climate exposures paediatric lung-function trajectories 2026}
{Climate change as a modifier of paediatric lung-function trajectories.}
{Eur Respir Rev}
{Mini-review using lung-function trajectories as its frame. Climate-sensitive exposures shift trajectories downward by three routes: early developmental disruption (gestational pollution and temperature, via growth restriction and preterm birth), sustained exposure impairing postnatal growth and peak, and exacerbations with incomplete recovery. Evidence that cleaner air accelerates lung-function growth is taken as proof the trajectory can be protected. No effect sizes or search method in the abstract. Limitation: co-occurring heat, smoke, allergen and damp are not separated, so attribution to PM is not possible from this frame. Relevance: low - frames why childhood exposure windows matter, without a measurement requirement.}
{10.1183/16000617.0222-2026}

\par\vskip 0pt plus 10\baselineskip\penalty-200\vskip 0pt plus -10\baselineskip
\band{Mechanistic toxicology\hfill 2 records}{Violet}

\paperentry{Li poultry-house PM$_{2.5}$-derived material APEC sequential exposure chick cell 2026}
{Inflammatory and redox responses associated with sequential exposure to poultry-house PM$_{2.5}$-derived material and avian pathogenic Escherichia coli in chicken models.}
{Poult Sci}
{Chicken type II-like lung cells and 40 SPF chicks (n=10/group, 14 days) received control, single or sequential PM$_{2.5}$-derived material then avian pathogenic E. coli. The sequential group showed the largest changes in several inflammatory and redox endpoints and lower CFTR; N-acetyl-L-cysteine and CBIQ shifted cell responses toward control. No effect sizes in the abstract. Limitation, stated by the authors: the eluted, lyophilized, autoclaved material was not characterized for size, composition or endotoxin, and schedules differed in timing, so no interaction, synergy or particle-specific effect is established. Relevance: low - a sterilized eluate is not an airborne particle.}
{10.1016/j.psj.2026.108064}

\paperentry{Qureshi narrative review bioaccessibility particle-bound metals PAHs tiered framework 2026}
{Bioaccessibility and Toxicity of Particle-Bound Metals and Polycyclic Aromatic Hydrocarbons: Sources, Respiratory Fate, Mechanisms, and Risk Assessment}
{Research Square (preprint)}
{Single-author critical narrative review, not peer-reviewed. Metal release varies with element, source, size, extraction conditions and solid-to-liquid ratio, and acidic lysosomal fluid usually dissolves more than near-neutral extracellular fluid; PAH release depends on hydrophobicity, soot binding, surfactants and proteins. Poorly soluble particles can stay active with little released. Proposes a tiered framework joining mass, source, size, speciation, compartment-specific bioaccessibility, oxidative potential and delivered dose. No search method or quantitative synthesis in the abstract, and differences in simulated lung fluids limit comparison. Relevance: low - explains why equal mass need not mean equal hazard, beyond what a mass sensor reports.}
{10.21203/rs.3.rs-11258579/v1}

\par\vskip 0pt plus 10\baselineskip\penalty-200\vskip 0pt plus -10\baselineskip
\band{Exposure assessment \& modelling\hfill 3 records}{Clay}

\paperentry{Wu penguin colony chlorophyll-a mediation PM$_{1}$ PM$_{2.5}$ PM$_{10}$ Ardley Island 2026}
{Phytoplankton-mediated coupling between penguin colonies and atmospheric particulate matter in Antarctic coastal waters.}
{Mar Pollut Bull}
{Synchronous ocean-atmosphere transects around Ardley Island. After controlling for spatial autocorrelation, chlorophyll-a remained associated with PM$_{2.5}$ (standardized coefficient 0.12). Across PM$_{1}$, PM$_{2.5}$ and PM$_{10}$ indirect NH4-Chl-a-PM effects exceeded direct ones, up to eightfold, but only at alpha 0.10, with 95\% CIs including zero. The authors nonetheless propose Chl-a as a proxy for the colony atmospheric footprint. Limitation: the mediation result does not exclude zero, the transects are one spatial sample, and particle composition is not reported, so marine biogenic origin is inferred. Relevance: low - a very clean background where optical PM readings are near any low-cost sensor's noise floor.}
{10.1016/j.marpolbul.2026.120512}

\paperentry{Kim CMAQ-UNIPAR secondary organic aerosol South Korea ASIA-AQ 2026}
{Prediction of Secondary Organic Aerosol Formation in South Korea using CMAQ-UNIPAR during the ASIA-AQ Campaign}
{Atmos Environ}
{\textsc{metadata only.} Metadata only; no abstract in Crossref at issue time, so no finding is asserted. The title places a CMAQ run with the UNIPAR SOA module against the ASIA-AQ campaign over South Korea. Carried because a mechanistic SOA module tested against an airborne and ground campaign bears on the organic fraction of PM$_{2.5}$, the component that most changes refractive index and hygroscopicity across a sensor's calibration range, and on the CTM surfaces health cohorts inherit. To be assessed when the Elsevier abstract deposits: bias against campaign OA and any change in simulated PM$_{2.5}$. Relevance: possible, unassessed.}
{10.1016/j.atmosenv.2026.122465}

\paperentry{Gantina remote sensing CO HC NOx soot heavy-duty trucks buses Indonesia 2026}
{Emission of CO, HC, NOx, and Soot from HDV Trucks and Buses using Remote Sensing in Indonesia}
{Aerosol Air Qual Res}
{\textsc{metadata only.} Metadata only; no abstract in Crossref at issue time, so no finding is asserted. The title reports on-road remote sensing of CO, HC, NOx and soot from heavy-duty trucks and buses in Indonesia. Carried because measured real-world fleet soot factors are scarce for Southeast Asia, and because they anchor inventory inputs where reference networks are thin and low-cost networks are the main expansion route. To be assessed when an abstract deposits: fleet size sampled, soot retrieval method and any emission-standard contrast. Relevance: possible, unassessed.}
{10.1007/s44408-026-00174-7}

\par\vskip 0pt plus 10\baselineskip\penalty-200\vskip 0pt plus -10\baselineskip
\band{Sensing, forecasting \& instrumentation\hfill 4 records}{Amber}

\paperentry{Soderblom Laplace-domain in situ network calibration Helsinki Antwerp SensEURCity 2026}
{Laplace domain dynamic calibration in city-wide sensor networks}
{Meas Sci Technol}
{A Laplace-domain method estimates pairwise dynamic relationships between sensor signals, separates a medium-related spatial effect, and solves an overdetermined calibration against one known network reference, with Monte Carlo uncertainty. Tested on Helsinki (linked to SMEAR III) and Antwerp (SensEURCity). Coefficients were obtained for temperature, RH and pressure in both networks, but PM calibration was feasible only on restricted sensor subsets because PM signals are localized and weakly correlated. Injected shifts could be flagged when geometry and correlations sufficed. No PM R2 or RMSE in the abstract, so the size of the PM failure cannot be graded. Relevance: direct - a metrology-chain calibration transfer fails first on the measurand networks are built for.}
{10.1088/1361-6501/aebdb6}

\paperentry{Kloppenborg ARMIE SPS30 node laboratory DustTrak SMPS APS source-out calibration 2026}
{Laboratory Evaluation of ARMIE, a Portable SPS30-Based Low-Cost Sensor Node for PM$_{2.5}$ Monitoring}
{Sensors}
{SPS30 nodes were co-located in the laboratory with a DustTrak and an SMPS+APS across candle, cooking, cigarette smoke and clean air, at regular and high humidity. Agreement with the DustTrak was r 0.93-0.98; with SMPS+APS it was moderate, r 0.74-0.94, with systematic underestimation at high concentration. Leave-one-sensor-out and leave-one-source-out calibration reduced mean error and narrowed limits of agreement. Limitation: a DustTrak is itself a photometer, so the high r partly reflects shared optical bias; the sizing reference is the honest comparator and shows the gap. No gravimetric or field arm, and r is not accuracy. Relevance: direct - leave-one-source-out is the right test for epidemiology nodes. \textsc{backfill}, Crossref 2 January 2026.}
{10.3390/s26010280}

\paperentry{Goyal PurpleAir PA-II kappa-Kohler Lorentz-Lorenz random forest moisture artefact 2026}
{A physics-aware digital twin for correcting moisture-induced material artifacts in PM$_{2.5}$ sensors}
{Discover Electronics}
{A pair of PA-II units beside one EPA FEM monitor for 150 days in an industrial area. kappa-Kohler growth and a Lorentz-Lorenz refractive index mixing rule motivate features for a polynomial-enhanced random forest. At RH 76.16\% optical intensity rose 59.5\%, summarised as a mass-intensity divergence factor of 4.33. Raw bias of +17.55\% fell to 0.13\% NMB, R2 0.9010. Limitation: two sensors at one site, single author, and the abstract gives no train/test separation by time, sensor or site, so 0.13\% cannot be read as field skill. Calling FEM 'gravimetric' is loose. Relevance: moderate - the physics is standard and testable; the transferability is untested. \textsc{backfill}, Crossref 26 September 2026.}
{10.1007/s44291-026-00286-9}

\paperentry{Mao wavelet attention LSTM indoor aerosol concentration prediction 2026}
{A Wavelet-Based Attention LSTM Framework for Aerosol Concentration Prediction in Indoor Environments}
{J Aerosol Sci}
{\textsc{metadata only.} Metadata only; no abstract in Crossref at issue time, so no finding is asserted. The title describes a wavelet decomposition feeding an attention LSTM to predict indoor aerosol concentration. Carried because short-horizon indoor prediction is one of the candidate edge-inference tasks for a sensor node, and J Aerosol Sci is one of the venues where such work is tied to a measured aerosol. What decides its value is unknown until the abstract deposits: the measuring instrument, whether validation is by held-out episode rather than random split, and the prediction horizon. Relevance: possible, unassessed.}
{10.1016/j.jaerosci.2026.106908}

\par\vskip 0pt plus 10\baselineskip\penalty-200\vskip 0pt plus -10\baselineskip
\band{Occupational \& indoor\hfill 4 records}{Amber}

\paperentry{Myers NIOSH CWHSP pneumoconiosis dusty job type underground coal 1991-2024}
{Assessment of pneumoconiosis by dusty job type among underground coal miners participating in the Coal Workers' Health Surveillance Program,1991-2024.}
{Front Public Health}
{30,319 underground miners in the national surveillance programme, 1991-2024, tenure-adjusted. Continuous miner operators had CWP PR 1.8 (1.6-2.0) and PMF PR 3.0 (2.3-4.0) against those never in the job; roof bolters 1.4 (1.3-1.6) and 1.4 (1.1-1.9). CWP prevalence was 8.4\% in Central Appalachia versus 1.5\% elsewhere, adjusted PR 3.9 (3.4-4.5); PMF PR 10 (5.7-17.5). Limitation: job title stands in for dust, and programme participants are self-selected, so healthy-worker and participation bias run in unknown directions. Prevalence pools 33 years of changing limits. Relevance: moderate - a tenfold regional contrast in a monitored regulated industry is the strongest argument today for personal rather than compliance-point dust data.}
{10.3389/fpubh.2026.1970741}

\paperentry{Molina Rueda AirPen facility-wide personal PM noise formaldehyde one shift 2026}
{Facility-wide assessment of occupational exposure to dust, noise, and formaldehyde}
{J Expo Sci Environ Epidemiol}
{Five people deployed 83 AirPen compact personal monitors across one full shift at a furniture plant of about 100 workers. Valid after screening: 67 PM, 67 noise, 22 formaldehyde (72, 67, 31 counting samples below the LOD). On-board sensors identified anomalies and confirmed wear. Distribution parameters were estimated with high certainty, and confidence intervals grew several-fold when fewer samples were analysed. Limitation: one facility on one shift says nothing about day-to-day variance, usually the main occupational term; 16 of 83 PM samples were lost or below LOD, and no comparison values are in the abstract. Relevance: direct - sample count, not instrument class, sets the uncertainty. \textsc{backfill}, entry 28 May 2026.}
{10.1038/s41370-026-00883-6}

\paperentry{Olsen LC-APCI-MS resin acids personal softwood sawmill dust Norway 2026}
{A Versatile Method for Determination of Resin Acids in Occupational Softwood Sawmill Dust Samples by Liquid Chromatography-Atmospheric Pressure Chemical Ionization-Mass Spectrometry.}
{J Sep Sci}
{Six reversed-phase columns and multiple mobile phases were screened; a C18-amide column with acetonitrile and 30 mM ammonium acetate at pH 6.30 and negative APCI gave a 15-minute isocratic run. Applied to six personal sawmill dust extracts, it determined abietic, dehydroabietic and levopimaric acids, coeluting isopimaric and pimaric acids, and 7-oxodehydroabietic acid. Limitation: six extracts demonstrate applicability, not exposure levels, and two isomers still coelute despite that being the problem the method targets. LODs are not in the abstract. Relevance: low - constituent chemistry of an occupational dust that no optical sensor could separate, useful only as a reference analysis.}
{10.1002/jssc.70568}

\paperentry{Apak office-to-education studio zone CO$_2$ PM$_{2.5}$ PM$_{10}$ TVOC perception Istanbul 2026}
{Indoor Air Quality and Occupant Perception in an Office Building Adapted for Educational Use: A Zone‐Based Case Study}
{Indoor Air}
{One mechanically ventilated design studio, three zones, questionnaire from 48 users. CO$_2$ was highest and most variable in the central zone (527 +/- 114 ppm, 95th percentile 760 ppm), and upper-end particulate indicators were greater in the circulation zone, but no zonal difference in CO$_2$, PM or TVOC survived multiplicity adjustment. Perceived freshness tracked overall satisfaction (Spearman rho 0.475, 0.219-0.682). Limitation, stated by the authors: respondents cannot be linked to measurement times or locations, so this is no exposure-perception relationship; sequential monitoring and few points restrict inference. Relevance: low - sequential single-instrument sampling cannot resolve within-room PM differences, which simultaneous nodes could.}
{10.1155/ina/4355214}

\clearpage
\section*{\sffamily\bfseries\color{Deep}Corpus register}
\vspace{-1.5mm}{\color{Amber}\rule{\linewidth}{1.1pt}}\vspace{2.5mm}

\input{table_rows}

\vspace{2mm}

\begin{keybox}
{\sffamily\bfseries\footnotesize\color{Deep}Provenance and method}\\[1.2mm]
{\fontsize{7.2}{8.8}\selectfont
Entry window \textbf{4 November 2026}, single day (Wednesday). Harvester legs: \textbf{PubMed} health \textbf{12}, sensing \textbf{0}; \textbf{Europe PMC} \textbf{17} (1 AGRICOLA retro-index record dropped); \textbf{Crossref by ISSN} \textbf{53}; \textbf{arXiv} no in-window entry; \textbf{OpenAlex} HTTP 429 on the harvester leg. 81 unique after cross-source dedupe, 8 already carried, \textbf{73 fresh}; PubMed connector (EDAT 4 Nov, three axes) 22 PMIDs, 21 bucketed to 4 Nov, 12 appended --- \textbf{85 screened}. Abstracts: 24 recovered by efetch; 20 records, mostly Elsevier day-0 deposits, stayed empty.
\textbf{Amass} BiomedCore (ingest-date filter) 15 returned: 7 already in \texttt{seen.json}, 1 previously rejected, \textbf{1 admitted as backfill}, 6 held for a later slot because the backfill cap was reached. \textbf{Consensus} calibration sweep: 10 hits, 1 new by DOI (ARMIE, admitted as backfill). The third backfill (Goyal) was held from the previous run's Consensus sweep. \textbf{OpenAlex, alphaXiv, bioRxiv/medRxiv and Wiley Scholar Gateway connectors were not available on this session}, so those legs did not run. \textbf{Scite} retraction check \textbf{unavailable} --- the connector again refused with a paid-plan error, so no record here carries a Scite retraction check.
\textbf{18 admitted} (15 window, 3 backfill), \textbf{70 rejected} with logged reasons. \textbf{Trial watch} 4--6 November: ClinicalTrials.gov 1 hit, NCT07914010 (REACT), already carried; Amass TrialCore 15 returned, relevant ones already carried --- \textbf{0 new}. \textbf{3 carried on metadata alone}. The x2 funnel is drawn (15 ratio estimates) but Egger is withheld: 15 estimates from 4 sources are not independent.
Abstracts remain the copyright of their publishers.\par}
\end{keybox}

\end{document}
